\documentclass{beamer}

\usepackage[T2A]{fontenc}
\usepackage[utf8]{inputenc}
\usepackage[english,russian]{babel}
\pdfmapfile{+sansmathaccent.map}


\mode<presentation>
{
	\usetheme{Warsaw} % or try Darmstadt, Madrid, Warsaw, Rochester, CambridgeUS, ...
	\usecolortheme{seahorse} % or try seahorse, beaver, crane, wolverine, ...
	\usefonttheme{serif}  % or try serif, structurebold, ...
	\setbeamertemplate{navigation symbols}{}
	\setbeamertemplate{caption}[numbered]
} 


%%%%%%%%%%%%%%%%%%%%%%%%%%%%
% itemize settings


%%%%%%%%%%%%%%%%%%%%%%%%%%%%
% itemize settings

\definecolor{myhotpink}{RGB}{255, 80, 200}
\definecolor{mywarmpink}{RGB}{255, 60, 160}
\definecolor{mylightpink}{RGB}{255, 80, 200}
\definecolor{mypink}{RGB}{255, 30, 80}
\definecolor{mydarkpink}{RGB}{155, 25, 60}

\definecolor{mypaleblue}{RGB}{240, 240, 255}
\definecolor{mylightblue}{RGB}{120, 150, 255}
\definecolor{myblue}{RGB}{90, 90, 255}
\definecolor{mygblue}{RGB}{70, 110, 240}
\definecolor{mycourseblue}{RGB}{120, 140, 255}
\definecolor{mycourselightblue}{RGB}{245, 250, 255}
\definecolor{mydarkblue}{RGB}{0, 0, 180}
\definecolor{myblackblue}{RGB}{40, 40, 120}

\definecolor{mygreen}{RGB}{0, 200, 0}
\definecolor{mygreen2}{RGB}{245, 255, 230}
\definecolor{mydarkgreen}{RGB}{0, 120, 0}
\definecolor{mydarkgreen}{RGB}{0, 120, 0}


\definecolor{mygray}{gray}{0.8}
\definecolor{mydarkgray}{RGB}{80, 80, 160}
\definecolor{mylightgray}{RGB}{160, 160, 160}

\definecolor{mydarkred}{RGB}{160, 30, 30}
\definecolor{mylightred}{RGB}{255, 150, 150}
\definecolor{myred}{RGB}{200, 110, 110}
\definecolor{myblackred}{RGB}{120, 40, 40}

\definecolor{mypink}{RGB}{255, 30, 80}
\definecolor{myhotpink}{RGB}{255, 80, 200}
\definecolor{mywarmpink}{RGB}{255, 60, 160}
\definecolor{mylightpink}{RGB}{255, 80, 200}
\definecolor{mydarkpink}{RGB}{155, 25, 60}
\definecolor{mywhitepink}{RGB}{255, 240, 240}

\definecolor{mydarkcolor}{RGB}{60, 25, 155}
\definecolor{mylightcolor}{RGB}{130, 180, 250}



\setbeamertemplate{itemize items}[default]

\setbeamertemplate{itemize item}{\color{myblackblue}$\blacksquare$}
\setbeamertemplate{itemize subitem}{\color{mygblue}$\blacktriangleright$}
\setbeamertemplate{itemize subsubitem}{\color{mygray}$\blacksquare$}

\setbeamercolor{palette quaternary}{fg=white,bg=mycourseblue}
\setbeamercolor{titlelike}{parent=palette quaternary}

\setbeamercolor{palette quaternary2}{fg=black,bg=mycourselightblue}
\setbeamercolor{frametitle}{parent=palette quaternary2}

\setbeamerfont{frametitle}{size=\Large,series=\scshape}
\setbeamerfont{framesubtitle}{size=\normalsize,series=\upshape}





%%%%%%%%%%%%%%%%%%%%%%%%%%%%
% block settings

\setbeamercolor{block title}{bg=red!30,fg=black}

\setbeamercolor*{block title example}{bg=mygreen!40!white,fg=black}

\setbeamercolor*{block body example}{fg= black, bg= mygreen2}


%%%%%%%%%%%%%%%%%%%%%%%%%%%%
% URL settings
\hypersetup{
	colorlinks=true,
	linkcolor=blue,
	filecolor=blue,      
	urlcolor=blue,
}

%%%%%%%%%%%%%%%%%%%%%%%%%%

\renewcommand{\familydefault}{\rmdefault}


\newcommand{\mycolor}[1] {
\setbeamercolor{normal text}{fg={#1}} % Recolors standard body text
\setbeamercolor{structure}{fg={#1}}   % Recolors bullets and structure items
\usebeamercolor[fg]{normal text}      % Applies the normal text color immediately
}


\usepackage{amsmath}
\usepackage{mathtools}

\usepackage{cancel}
\usepackage{centernot}

\usepackage{changepage}
\usepackage{subcaption}

\usepackage{qrcode}

\newcommand{\bo}[1] {\mathbf{#1}}
\newcommand{\R}{\mathbb{R}} 
\newcommand{\T}{^\top}     



\newcommand{\mydate}{Fall 2026}

\newcommand{\mygit}{\textcolor{blue}{\href{https://github.com/SergeiSa/Advanced-Control-2026}{github.com/SergeiSa/Advanced-Control-2026}}}

\newcommand{\myqr}{ \textcolor{black}{\qrcode[height=1.5in]{https://github.com/SergeiSa/Advanced-Control-2026}}
}




\newcommand{\bref}[2] {\textcolor{blue}{\href{#1}{#2}}}

%%%%%%%%%%%%%%%%%%%%%%%%%%%%
% code settings

\usepackage{listings}
\usepackage{color}
% \definecolor{mygreen}{rgb}{0,0.6,0}
% \definecolor{mygray}{rgb}{0.5,0.5,0.5}
\definecolor{mymauve}{rgb}{0.58,0,0.82}
\lstset{ 
	backgroundcolor=\color{white},   % choose the background color; you must add \usepackage{color} or \usepackage{xcolor}; should come as last argument
	basicstyle=\footnotesize,        % the size of the fonts that are used for the code
	breakatwhitespace=false,         % sets if automatic breaks should only happen at whitespace
	breaklines=true,                 % sets automatic line breaking
	captionpos=b,                    % sets the caption-position to bottom
	commentstyle=\color{mygreen},    % comment style
	deletekeywords={...},            % if you want to delete keywords from the given language
	escapeinside={\%*}{*)},          % if you want to add LaTeX within your code
	extendedchars=true,              % lets you use non-ASCII characters; for 8-bits encodings only, does not work with UTF-8
	firstnumber=0000,                % start line enumeration with line 0000
	frame=single,	                   % adds a frame around the code
	keepspaces=true,                 % keeps spaces in text, useful for keeping indentation of code (possibly needs columns=flexible)
	keywordstyle=\color{blue},       % keyword style
	language=Octave,                 % the language of the code
	morekeywords={*,...},            % if you want to add more keywords to the set
	numbers=left,                    % where to put the line-numbers; possible values are (none, left, right)
	numbersep=5pt,                   % how far the line-numbers are from the code
	numberstyle=\tiny\color{mygray}, % the style that is used for the line-numbers
	rulecolor=\color{black},         % if not set, the frame-color may be changed on line-breaks within not-black text (e.g. comments (green here))
	showspaces=false,                % show spaces everywhere adding particular underscores; it overrides 'showstringspaces'
	showstringspaces=false,          % underline spaces within strings only
	showtabs=false,                  % show tabs within strings adding particular underscores
	stepnumber=2,                    % the step between two line-numbers. If it's 1, each line will be numbered
	stringstyle=\color{mymauve},     % string literal style
	tabsize=2,	                   % sets default tabsize to 2 spaces
	title=\lstname                   % show the filename of files included with \lstinputlisting; also try caption instead of title
}


%%%%%%%%%%%%%%%%%%%%%%%%%%%%
% URL settings
\hypersetup{
	colorlinks=false,
	linkcolor=blue,
	filecolor=blue,      
	urlcolor=blue,
}

%%%%%%%%%%%%%%%%%%%%%%%%%%

%%%%%%%%%%%%%%%%%%%%%%%%%%%%
% tikz settings

\usepackage{tikz}
\tikzset{every picture/.style={line width=0.75pt}}
\usetikzlibrary{shapes.geometric, decorations.markings, positioning, patterns, calc, arrows.meta}

\newcommand{\myqrframe}{
	\begin{frame}
		\begin{center}
			\centerline{Слайды доступны на Github:}
			\bigskip
			\centerline{\mygit}
			\bigskip
			\myqr
		\end{center}
	\end{frame}
}


\subtitle{Современные методы управления роботами}
\author{Сергей Савин}
\date{2026}

\title{Метод сопряжённых переменных, распространение чувствительностей}




\begin{document}
	\maketitle
	
	
	\begin{frame}{Содержание}
		\textbf{Метод сопряжённых переменных}
		\begin{itemize}
			\item Вывод через лагранжиан и гамильтониан
			\item Уравнения для костояния и итоговый градиент
			\item Сравнение с принципом Понтрягина
			\item Связь с обратным режимом AD
		\end{itemize}
		
		\textbf{Метод распространения чувствительностей}
		\begin{itemize}
			\item Вывод ОДУ для чувствительностей
			\item Вычисление градиента стоимости
			\item Связь с прямым режимом AD
		\end{itemize}
	\end{frame}
	
	
	
	%------------------------------------------------
	\begin{frame}{Прямое пристреливание: постановка задачи}
		\begin{flushleft}
			Планирование траектории методом прямого пристреливания можно представить как
			следующую задачу оптимизации:
			
			\[
			\boxed{
				\begin{aligned}
					\min_{c}\quad
					J(c)&=\Phi(x(t_f))+
					\int_0^{t_f}\ell(x,u,t)\,dt\\
					\text{s.t.}\quad
					\dot{x}&=f(x,u,t),\\
					u&=u(c,t),\\
					x(0)&=x_0.
				\end{aligned}
			}
			\]
			
			\bigskip
			
			Здесь $c$ параметризует управляющее воздействие, а
			$x(t)$ определяется интегрированием динамики. $\Phi$ --- терминальная стоимость, $\ell$ --- интегральная стоимость.
		\end{flushleft}
	\end{frame}
	
	%------------------------------------------------
	\begin{frame}{Градиенты}
		\begin{flushleft}
			Чтобы найти параметры $c$, минимизирующие $J(c)$,
			нам нужно вычислить градиент $\frac{dJ}{dc}$.
			
			\bigskip
			
			Вычислить его явно сложно, потому что:
			\begin{itemize}
				\item Состояние $x(t)$ неявно зависит от $c$
				через динамику системы.
				\item Вычисление $dx(t)/dc$ требует дифференцирования
				решения ОДУ по каждому параметру.
				\item Число уравнений для чувствительностей растёт
				с числом параметров в $c$.
			\end{itemize}
			
			\bigskip
			
			Мы выведем метод, который позволяет избежать явного
			вычисления $dx(t)/dc$.
		\end{flushleft}
	\end{frame}
	
	%------------------------------------------------
	\begin{frame}{Градиенты через лагранжиан}
		\begin{flushleft}
			Введём лагранжиан, чтобы учесть ограничение-равенство,
			описывающее динамику системы:
			
			\begin{align}
				\boxed{
					\begin{aligned}
						\mathcal{L}(c,\lambda)
						={}&\Phi(x(t_f))\\
						&+\int_0^{t_f}
						\left[
						\ell(x,u,t)+
						\lambda^\top(f(x,u,t)-\dot{x})
						\right]dt.
					\end{aligned}
				}
			\end{align}
			
			\medskip
			
			Здесь $\lambda(t)$ --- зависящий от времени множитель Лагранжа,
			также называемый \textbf{костоянием}.
			
			\bigskip
			
			На множестве $\dot{x}=f(x,u,t)$, которое включает все
			допустимые траектории,
			\[
			\boxed{\mathcal{L}(c,\lambda)=J(c).}
			\]
		\end{flushleft}
	\end{frame}
	
	%------------------------------------------------
	\begin{frame}{Упрощение лагранжиана}
		\begin{flushleft}
			Ключевой шаг --- исключить $\dot{x}$ с помощью интегрирования
			по частям.
			
			\begin{theorem}[Интегрирование по частям]
				\[
				\int_a^b u^\top\dot{v}\,dt
				=
				\left[u^\top v\right]_a^b
				-\int_a^b \dot{u}^\top v\,dt.
				\]
			\end{theorem}
			
			Применение этой формулы даёт
			\[
			-\int_0^{t_f}\lambda^\top\dot{x}\,dt
			=
			-\lambda(t_f)^\top x(t_f)
			+\lambda(0)^\top x_0
			+\int_0^{t_f}\dot{\lambda}^\top x\,dt.
			\]
			
			Таким образом, лагранжиан принимает вид
			\[
			\boxed{
				\begin{aligned}
					\mathcal{L}
					={}&\Phi(x(t_f))-\lambda(t_f)^\top x(t_f)+\lambda(0)^\top x_0\\
					&+\int_0^{t_f}
					\left[\ell+\lambda^\top f+
					\dot{\lambda}^\top x\right]dt.
				\end{aligned}
			}
			\]
		\end{flushleft}
	\end{frame}
	
	%------------------------------------------------
	\begin{frame}{Градиент по параметрам}
		\begin{flushleft}
			Продифференцируем $\mathcal{L}$ по $c$,
			считая $\lambda(t)$ фиксированным:
			
			\[
			\begin{aligned}
				\frac{dJ}{dc}
				={}&
				\left(\frac{d\Phi}{dx}
				-\lambda(t_f)^\top\right)
				\frac{dx(t_f)}{dc}\\
				&+\lambda(0)^\top\frac{dx_0}{dc}\\
				&+\int_0^{t_f}
				\left(
				\frac{d\ell}{dx}
				+\lambda^\top\frac{df}{dx}
				+\dot{\lambda}^\top
				\right)\frac{dx}{dc}\,dt\\
				&+\int_0^{t_f}
				\left(
				\frac{d\ell}{du}
				+\lambda^\top\frac{df}{du}
				\right)\frac{du}{dc}\,dt.
			\end{aligned}
			\]
			
			\bigskip
			
			Поскольку $x(0)=x_0$ не зависит от $c$,
			\[
			\frac{dx_0}{dc}=0.
			\]
		\end{flushleft}
	\end{frame}
	
	%------------------------------------------------
	\begin{frame}{Упрощение через гамильтониан}
		\begin{flushleft}
			Определим гамильтониан:
			\[
			\boxed{
				H(x,u,\lambda,t)=\ell(x,u,t)+\lambda^\top f(x,u,t).
			}
			\]
			
			Его производные равны
			\[
			\begin{aligned}
				\frac{dH}{dx}
				&=\frac{d\ell}{dx}+\lambda^\top\frac{df}{dx},\\
				\frac{dH}{du}
				&=\frac{d\ell}{du}+\lambda^\top\frac{df}{du}.
			\end{aligned}
			\]
			
			\bigskip
			
			Теперь градиент можно записать как
			\[
			\boxed{
				\begin{aligned}
					\frac{dJ}{dc}
					={}&
					\left(\frac{d\Phi}{dx}-\lambda(t_f)^\top\right)
					\frac{dx(t_f)}{dc}\\
					&+\int_0^{t_f}
					\left(\frac{dH}{dx}+\dot{\lambda}^\top\right)
					\frac{dx}{dc}\,dt\\
					&+\int_0^{t_f}\frac{dH}{du}\frac{du}{dc}\,dt.
				\end{aligned}
			}
			\]
		\end{flushleft}
	\end{frame}
	
	%------------------------------------------------
	\begin{frame}{Исключение градиентов состояния}
		\begin{flushleft}
			Выберем $\lambda(t)$ так, чтобы сократить все члены, содержащие
			$dx/dc$.
			
			\bigskip
			
			Сокращение терминального члена даёт
			\[
			\boxed{
				\lambda(t_f)=\left(\frac{d\Phi}{dx}\right)^\top.
			}
			\]
			
			\bigskip
			
			Сокращение интегрального члена даёт
			\[
			\boxed{
				\dot{\lambda}
				=-\left(\frac{dH}{dx}\right)^\top.
			}
			\]
			
			\bigskip
			
			Это ОДУ, решаемое \textbf{в обратном времени},
			с $\lambda(t_f)$ в качестве терминального условия.
		\end{flushleft}
	\end{frame}
	
	%------------------------------------------------
	\begin{frame}{Градиенты: итоговая форма}
		\begin{flushleft}
			При таком выборе костояния все члены, содержащие
			$dx/dc$, обращаются в нуль.
			
			\bigskip
			
			Получаем итоговое выражение для градиента:
			\[
			\boxed{
				\frac{dJ}{dc}
				=
				\int_0^{t_f}
				\frac{dH}{du}\frac{du}{dc}\,dt.
			}
			\]
			
			\bigskip
			
			Эквивалентно,
			\[
			\boxed{
				\frac{dJ}{dc}
				=
				\int_0^{t_f}
				\left(
				\frac{d\ell}{du}
				+\lambda^\top\frac{df}{du}
				\right)
				\frac{du}{dc}\,dt.
			}
			\]
			
			Чувствительности состояния $dx/dc$ больше не нужны.
		\end{flushleft}
	\end{frame}
	
	%------------------------------------------------
	\begin{frame}{Описание метода сопряжённых переменных}
		\begin{flushleft}
			Метод сопряжённых переменных вычисляет градиент за четыре шага:
			
			\begin{enumerate}
				\item \textbf{Прямой ход:}
				Интегрируем $\dot{x}=f(x,u,t)$ и сохраняем траекторию.
				
				\item \textbf{Обратный ход:}
				Интегрируем $\dot{\lambda}=-(dH/dx)^\top$ в обратном времени,
				начиная с $\lambda(t_f)=(d\Phi/dx)^\top$.
				
				\item \textbf{Вычисление градиента:}
				Вычисляем $\displaystyle
				\frac{dJ}{dc}=\int_0^{t_f}
				\frac{dH}{du}\frac{du}{dc}\,dt$.
				
				\item \textbf{Обновление решателя:}
				Обновляем $c$ с помощью градиентного оптимизатора
				и повторяем до сходимости.
			\end{enumerate}
			
			\bigskip
			
			Стоимость обратного прохода не растёт с
			числом параметров управления.
		\end{flushleft}
	\end{frame}
	
	%------------------------------------------------
	\begin{frame}{Сравнение с принципом Понтрягина}
		\begin{flushleft}
			
			\textbf{Принцип минимума Понтрягина} даёт необходимые
			условия оптимальности траектории управления. Для оптимального решения он вводит гамильтониан
			и костояние, удовлетворяющие
			\[
			\dot{\lambda}=-\left(\frac{dH}{dx}\right)^\top,
			\qquad
			\lambda(t_f)=\left(\frac{d\Phi}{dx}\right)^\top.
			\]
			
			Для безусловного оптимального управления он также требует
			\[
			\frac{dH}{du}=0
			\]
			вдоль оптимальной траектории.
			
			\textbf{Ключевое различие:}
			\begin{itemize}
				\item Метод сопряжённых переменных вычисляет градиент
				$dJ/dc$ для выбранного параметризованного управления.
				\item Принцип Понтрягина даёт необходимые
				условия оптимальности для самого оптимального управления.
			\end{itemize}
			
		\end{flushleft}
	\end{frame}
	%------------------------------------------------
	
	
	
	%------------------------------------------------
	\begin{frame}{Обратный режим автоматического дифференцирования vs.\ метод сопряжённых переменных}
		\begin{flushleft}
			
			\textbf{Дискретный случай (пристреливание с транскрипцией):} Обратный режим AD на развёрнутом цикле интегрирования
			\emph{и есть} дискретный метод сопряжённых переменных.
			
			\bigskip
			
			\textbf{Ограничения реализации AD:} Обратный режим AD строит вычислительный граф (часто большой), для хранения промежуточных значений может потребоваться значительный объём памяти.
			
		\end{flushleft}
	\end{frame}
	%------------------------------------------------
	
	
	
	
	%------------------------------------------------
	\begin{frame}{Прямое пристреливание: чувствительность}
		\begin{flushleft}
			Траектория состояния зависит от параметров управления
			через динамику и вход $u(c,t)$.
			
			Для каждого параметра $c_j$ определим чувствительность состояния
			\[
			S_j(t) := \frac{dx(t)}{dc_j}.
			\]
			
			\bigskip
			
			\textbf{Идея:} Вывести ОДУ, описывающее эволюцию
			$S_j(t)$. Как только чувствительности известны, можно вычислить
			градиент $\frac{dJ}{dc}$ по правилу цепочки.
		\end{flushleft}
	\end{frame}
	
	%------------------------------------------------
	\begin{frame}{Прямое пристреливание: дифференцирование динамики}
		\begin{flushleft}
			Начнём с динамики
			\[
			\dot{x} = f(x,u,t),
			\qquad u=u(c,t).
			\]
			
			Продифференцируем обе стороны по $c_j$:
			\[
			\frac{d}{dc_j}\dot{x}
			=
			\frac{df}{dx}\frac{dx}{dc_j}
			+
			\frac{df}{du}\frac{du}{dc_j}.
			\]
			
			Поскольку $c_j$ не зависит от времени, дифференцирование
			по $c_j$ коммутирует с дифференцированием
			по $t$:
			\[
			\frac{d}{dc_j}\dot{x}
			=
			\frac{d}{dt}\frac{dx}{dc_j}.
			\]
			
			Определяя $S_j := \frac{dx}{dc_j},
			\quad
			U_j := \frac{du}{dc_j}$, 
			получаем:
			\[
			\boxed{
				\dot{S}_j
				=
				\frac{df}{dx}S_j
				+
				\frac{df}{du}U_j
			}
			\]
		\end{flushleft}
	\end{frame}
	
	%------------------------------------------------
	\begin{frame}{Прямое пристреливание: ОДУ для чувствительностей}
		\begin{flushleft}
			Соберём чувствительности по всем $N$
			параметрам управления в матрицы:
			\[
			S =
			\begin{bmatrix}
				S_1 & \cdots & S_N
			\end{bmatrix}
			=
			\frac{dx}{dc},
			\qquad
			U_c = \frac{du}{dc}.
			\]
			
			Объединение отдельных уравнений для чувствительностей даёт
			\[
			\boxed{
				\dot{S}(t)
				=
				\frac{df}{dx}S(t)
				+
				\frac{df}{du}U_c(t)
			}
			\]
			
			\bigskip
			
			Поскольку $x(0)=x_0$ не зависит от $c$,
			\[
			\boxed{S(0)=0.}
			\]
			
			\textbf{Численная процедура:} Интегрируем ОДУ состояния и
			чувствительностей совместно в прямом времени. Вычисляем $\frac{df}{dx}$
			и $\frac{df}{du}$ вдоль траекторий состояния и входа.
		\end{flushleft}
	\end{frame}
	
	%------------------------------------------------
	\begin{frame}{Прямое пристреливание: градиент стоимости}
		\begin{flushleft}
			Целевая функция имеет вид
			\[
			J(c)=\Phi(x(t_f))
			+\int_0^{t_f}\ell(x,u,t)\,dt.
			\]
			
			Применим правило цепочки, используя
			$S(t)=\frac{dx}{dc}$ и $U_c(t)=\frac{du}{dc}$:
			\[
			\boxed{
				\begin{aligned}
					\frac{dJ}{dc}
					={}&
					\frac{d\Phi}{dx}S(t_f)\\
					&+
					\int_0^{t_f}
					\left(
					\frac{d\ell}{dx}S(t)
					+
					\frac{d\ell}{du}U_c(t)
					\right)dt.
				\end{aligned}
			}
			\]
			
			\bigskip
			
			Таким образом, ОДУ состояния и чувствительностей дают
			траекторию и её производные, из которых мы вычисляем
			$\frac{dJ}{dc}$.
		\end{flushleft}
	\end{frame}
	
	
	
	%------------------------------------------------
	\begin{frame}{Распространение чувствительностей vs прямой режим AD}
		\begin{flushleft}
			
			Распространение чувствительностей в дискретном времени --- это \emph{прямой режим автоматического дифференцирования}:
			
			\bigskip
			
			Для дискретной системы $x_{k+1}=F_k(x_k,c)$ и распространение чувствительностей, и прямой режим AD представляются как:
			
			\[ \boxed{ S_{k+1} = \frac{dF_k}{dx_k}S_k + \frac{dF_k}{dc}, \qquad S_0=0. } \]
			
			
		\end{flushleft}
	\end{frame}
	
	
	
	
	
	\begin{frame}{Литература}
		
		\begin{itemize}
			
			
			\item Russ Tedrake -- MIT:
			\bref{https://underactuated.mit.edu/trajopt.html}
			{Underactuated Robotics: Trajectory Optimization, ch. 10.}
			
		\end{itemize}
		
	\end{frame}
	
	
	

\myqrframe




\end{document}
